%% REVTeX 4-2 — Physical Review E
%% TCSV Paper II — Hierarchical scaling and epilepsy
%% Compile: pdflatex → bibtex → pdflatex × 2
\documentclass[aps,pre,twocolumn,showpacs,superscriptaddress,longbibliography]{revtex4-2}

\usepackage{amsmath,amssymb,amsfonts}
\usepackage{graphicx}
\usepackage{dcolumn}
\usepackage{bm}
\usepackage{hyperref}
\usepackage{xcolor}
\usepackage{booktabs}
\usepackage{multirow}

%% ── Custom commands ─────────────────────────────────────────
\newcommand{\kN}{k/N}
\newcommand{\taupop}{\tau_{\mathrm{pop}}}
\newcommand{\tauISI}{\tau_{\mathrm{ISI}}}
\newcommand{\tauabs}{\tau_{\mathrm{abs}}}
\newcommand{\hneu}{h_{\mathrm{neural}}}
\newcommand{\fUV}{f_{\mathrm{UV}}}
\newcommand{\rhoinfo}{\rho_{\mathrm{info}}}
\newcommand{\lp}{\lambda_{+}}

\begin{document}

%% ════════════════════════════════════════════════════════════
%%  TITLE & AUTHORS
%% ════════════════════════════════════════════════════════════
\title{Hierarchical scaling of cortical metabolic constraints
       predicts multi-scale neural oscillations and
       the energetic collapse in epilepsy}

\author{Gesivaldo Santos}
\email{gesi@uesb.edu.br}
\affiliation{Department of Biological Science,
             State University of Southwest Bahia (UESB),
             Jequié, Bahia, Brazil}

\date{\today}

%% ════════════════════════════════════════════════════════════
%%  ABSTRACT
%% ════════════════════════════════════════════════════════════
\begin{abstract}
In a companion paper~\cite{Santos2025a}, we showed that the
spectral structure of synaptic connectivity matrices follows
a universal power law $\lambda_k \sim k^{-0.43}$, conserved
across cortical species and areas, arising from the metabolic
constraint $\kN \approx 0.15$ that limits simultaneous neuronal
firing in the cortical microcircuit. Here we show that this
constraint scales hierarchically across all spatial and temporal
scales of cortical organization --- from individual microcircuits
($N \approx 100$ neurons, $\tauISI \approx 50\,\mathrm{ms}$)
through cortical areas ($N \approx 10^4$ minicolumns,
$\tau \approx 150\,\mathrm{ms}$) to the entire cortex
($N \approx 10^5$ regions, $\tau \approx 5{,}000\,\mathrm{ms}$).
The same quantization equation $f_n = 1/(n\taupop)$,
$n = 1,\ldots,6$, with $\taupop = \tauISI\sqrt{\kN}$,
predicts gamma oscillations ($52\,\mathrm{Hz}$) at the
microcircuit scale and alpha/beta oscillations ($3$--$17\,\mathrm{Hz}$)
at the cortical area scale. The constraint $\kN = 0.15$ is
invariant across all scales, as it is determined by the same
metabolic budget of cortical tissue. Critically, violation of
this constraint at any scale leads to energetic collapse: if
$\kN \to 1$, the required metabolic rate increases $6.7\times$
above the resting baseline, exceeding the maximum cerebrovascular
reserve ($3$--$4\times$) and producing tissue hypoxia within
$10$--$15\,\mathrm{s}$. We identify this as the thermodynamic
mechanism of epileptic seizure termination: the tonic-clonic
seizure is not merely a disorder of excessive neuronal firing,
but a violation of the metabolic UV constraint at progressively
larger spatial scales, culminating in whole-cortex collapse and
ictal silence. The semantic inertia $\hneu$ scales from
$5.47 \times 10^{-14}\,\mathrm{J\cdot s}$ at the microcircuit
to $1.51 \times 10^{-9}\,\mathrm{J\cdot s}$ at the whole-cortex
scale --- a $27{,}600\times$ amplification that formalizes the
thermodynamic resistance of large-scale cortical representations
to perturbation. These results establish $\kN = 0.15$ as a
multi-scale invariant of cortical dynamics, with the spectral
power law $\lambda_k \sim k^{-0.43}$ as its eigenvalue
signature at every hierarchical level.
\end{abstract}

\pacs{87.19.lj, 87.19.xj, 87.19.le, 87.10.Mn}
\keywords{cortical hierarchy, metabolic constraint, epilepsy,
scaling, UV catastrophe, sparse coding, neural thermodynamics,
semantic inertia}

\maketitle

%% ════════════════════════════════════════════════════════════
%%  I. INTRODUCTION
%% ════════════════════════════════════════════════════════════
\section{\label{sec:intro}Introduction}

\subsection{The metabolic constraint as a fundamental invariant}

The neocortical microcircuit operates under a strict metabolic
budget derived from the energetics of neural signaling.
Attwell and Laughlin~\cite{Attwell2001} established that the
dominant energy costs of cortical grey matter are allocated as:
$c_T = 0.47$ (action potentials), $c_V = 0.34$ (synaptic
transmission), $\mu = 0.13$ (resting potential maintenance),
and $\gamma = 0.03$ (other processes). From these coefficients,
the fraction of neurons that can simultaneously fire within the
aerobic metabolic budget is $\kN \approx 0.15$.

In a companion paper~\cite{Santos2025a}, we showed that this
constraint generates a discrete set of collective synaptic
modes (attractors) whose eigenvalue spectrum follows a universal
power law $\lambda_k \sim k^\alpha$ with $\alpha \approx -0.43$,
conserved between Blue Brain Project simulations and MICrONS
electron microscopy data. The constraint also quantizes the
cortical oscillation spectrum into six canonical EEG frequency
bands via $f_n = 1/(n\taupop)$, $n = 1, \ldots, 6$, where
$\taupop = \tauISI\sqrt{\kN}$ is the population integration time.

\subsection{Epilepsy as the central motivation}

Epilepsy affects approximately 50 million people
worldwide~\cite{WHO2023} and remains poorly understood at
the mechanistic level. Current models describe seizures
primarily in terms of E/I imbalance, hypersynchrony, or ion
channel dysfunction~\cite{Stafstrom2010}, but none provides a
quantitative thermodynamic account of why seizures terminate
spontaneously and why they produce the characteristic sequence
of tonic activity, clonic activity, and postictal silence.

We propose that epileptic seizure is the macroscopic
manifestation of the metabolic UV catastrophe of the cortical
ensemble: the progressive violation of $\kN \leq 0.15$ at
increasing spatial scales, until the whole-cortex metabolic
reserve is exhausted and the system undergoes forced quiescence.
This framework makes four quantitative predictions:
\begin{enumerate}
\item Seizure duration before forced termination:
  $t_{\mathrm{ictus}} \approx 10$--$15\,\mathrm{s}$ per scale.
\item Metabolic amplification at full activation:
  $1/\kN_{\mathrm{rest}} = 6.7\times$, exceeding the
  cerebrovascular reserve ($3$--$4\times$).
\item The spectral power law $\lambda_k \sim k^{-0.43}$
  should steepen to $\alpha \to -0.53$ in ictal LFP
  immediately before seizure termination.
\item Postictal silence duration scales with the spatial
  extent of the seizure.
\end{enumerate}

\subsection{Semantic inertia and resistance to change}

The semantic inertia $\hneu = \langle E \rangle \taupop / H$,
where $H$ is the activation entropy of the ensemble, quantifies
the thermodynamic cost of changing a neural representation.
At the microcircuit scale, $\hneu \approx 5.47 \times
10^{-14}\,\mathrm{J\cdot s}$. At the whole-cortex scale,
$\hneu \approx 1.51 \times 10^{-9}\,\mathrm{J\cdot s}$ ---
a $27{,}600\times$ amplification that formalizes the resistance
of large-scale cortical representations to
perturbation~\cite{McClelland1995}.

%% ════════════════════════════════════════════════════════════
%%  II. THEORETICAL FRAMEWORK
%% ════════════════════════════════════════════════════════════
\section{\label{sec:theory}Theoretical Framework}

\subsection{The TCSV thermodynamic ensemble}

The cortical microcircuit of $N$ neurons is modeled as a
canonical ensemble with metabolic budget constraint.
The Lagrangian of the system is~\cite{Santos2025a}:
\begin{equation}
  L(q,\dot{q}) = \frac{1}{2}c_T\rhoinfo\dot{q}^2
               - \frac{1}{2}c_V k q^2,
  \label{eq:lagrangian}
\end{equation}
where $q$ is the generalized cognitive coordinate and
$\rhoinfo = V_{\mathrm{syn}}/\langle V_{\mathrm{syn}}\rangle$
is the informational density.
The population integration time and quantized modes are:
\begin{equation}
  \taupop = \tauISI \cdot \sqrt{\kN},
  \label{eq:taupop}
\end{equation}
\begin{equation}
  f_n = \frac{1}{n \cdot \taupop}, \quad n = 1,\ldots,
  n_{\max} = \left\lfloor \frac{N}{k} \right\rfloor = 6.
  \label{eq:modes}
\end{equation}

\subsection{Hierarchical extension}

At each level of the cortical hierarchy
(Fig.~\ref{fig:hierarchy}), the quantized modes are:
\begin{equation}
  f_n^{(s)} = \frac{1}{n \cdot \tauISI^{(s)} \cdot \sqrt{\kN}},
  \quad n = 1,\ldots,6,
  \label{eq:modes_scaled}
\end{equation}
where $\kN = 0.15$ is scale-invariant and $\tauISI^{(s)}$ is
the characteristic inter-unit communication time at scale $s$.
The key hypothesis is that $\kN$ is scale-invariant, as the
aerobic metabolic budget applies at every spatial scale of
cortical tissue. The resulting hierarchy spans four orders of
magnitude in frequency (Table~\ref{tab:hierarchy}).

\begin{table}[h]
\caption{\label{tab:hierarchy}Quantized oscillation modes
across cortical scales.}
\begin{ruledtabular}
\begin{tabular}{lrrr}
Scale & $\tauISI^{(s)}$ & $f_1^{(s)}$ & $f_6^{(s)}$ \\
\hline
Microcircuit    &  49.8\,ms &  51.8\,Hz &  8.6\,Hz \\
Cortical area   & 150\,ms   &  17.2\,Hz &  2.9\,Hz \\
Cortical region & 800\,ms   &   3.2\,Hz &  0.5\,Hz \\
Whole cortex    & 5000\,ms  &   0.52\,Hz & 0.086\,Hz \\
\end{tabular}
\end{ruledtabular}
\end{table}

\subsection{The metabolic UV catastrophe}

The metabolic cost scales with activation fraction as:
\begin{equation}
  P(\kN) = P_{\mathrm{rest}} \cdot \frac{\kN}{0.15}.
  \label{eq:power}
\end{equation}
Full activation ($\kN = 1$) requires $6.7\times$ the resting
rate ($12{,}000\,\mathrm{mW}$), far exceeding the maximum
cerebrovascular reserve of $\sim 3$--$4\times$~\cite{Buzsaki2007},
producing hypoxia within $10$--$15\,\mathrm{s}$.

Two distinct spectral cutoffs emerge at each scale: an
\emph{ensemble cutoff} (Debye analog, state-dependent):
\begin{equation}
  \fUV^{\mathrm{ens}} = \frac{N/k - 1}{\tauISI},
  \label{eq:fUV_ens}
\end{equation}
and a \emph{neuronal absolute cutoff} (Planck analog, invariant):
\begin{equation}
  \fUV^{\mathrm{neural}} = \frac{1}{\tauabs\sqrt{\kN}}
  \approx 587\,\mathrm{Hz},
  \label{eq:fUV_neural}
\end{equation}
determined by the absolute refractory period
$\tauabs \approx 4.4\,\mathrm{ms}$.

\subsection{Semantic inertia at macroscopic scales}

The semantic inertia at scale $s$ is:
\begin{equation}
  \hneu^{(s)} = \frac{\langle E^{(s)}\rangle \cdot \taupop^{(s)}}
                      {H^{(s)}},
  \label{eq:h_scale}
\end{equation}
where $H^{(s)} = \ln\binom{N_s}{k_s}$ is the activation entropy.
For large $N_s$:
\begin{equation}
  \hneu^{(s)} \propto N_s \cdot \tauISI^{(s)} \cdot \sqrt{\kN},
  \label{eq:h_asymp}
\end{equation}
yielding the $27{,}600\times$ amplification from microcircuit
to whole cortex (Table~\ref{tab:h_scale}, Fig.~\ref{fig:h_neural}).

%% ════════════════════════════════════════════════════════════
%%  III. EPILEPSY AS UV CATASTROPHE
%% ════════════════════════════════════════════════════════════
\section{\label{sec:epilepsy}Epilepsy as the Metabolic
         UV Catastrophe}

\subsection{Phenomenology of generalized tonic-clonic seizures}

A generalized tonic-clonic seizure (GTCS) proceeds through:
(1)~\emph{tonic phase} ($\sim 10$--$20\,\mathrm{s}$), sustained
high-frequency firing ($>100\,\mathrm{Hz}$) across large cortical
areas; (2)~\emph{clonic phase} ($\sim 30$--$60\,\mathrm{s}$),
rhythmic burst-suppression at progressively lower frequencies;
and (3)~\emph{postictal silence} (minutes to hours), global
cortical suppression with gradual recovery~\cite{Krishnan2011}.
Current models attribute termination to synaptic depression or
active inhibition but provide no quantitative thermodynamic
account of the timing.

\subsection{The cascade of UV catastrophes
            (Fig.~\ref{fig:epilepsy})}

\textbf{Phase 1 --- Focal initiation ($\kN \to 0.5$).}
The ensemble UV cutoff drops from $114\,\mathrm{Hz}$ at rest to:
\begin{equation}
  \fUV^{\mathrm{ens}}(\kN=0.5) = \frac{1/0.5 - 1}{49.8\,\mathrm{ms}}
  \approx 20\,\mathrm{Hz}.
\end{equation}
Gamma suppression followed by low-frequency activity is the
ictal onset pattern observed in stereo-EEG.

\textbf{Phase 2 --- Regional propagation ($\kN \to 0.9$).}
\begin{equation}
  \fUV^{\mathrm{ens}}(\kN=0.9) \approx 2.2\,\mathrm{Hz}.
\end{equation}
Only delta activity is sustainable --- the clinical tonic phase.

\textbf{Phase 3 --- Whole-cortex collapse ($\kN \to 1$).}
Metabolic demand reaches $6.7\times$ the cerebrovascular
reserve. Tissue hypoxia develops within $10$--$15\,\mathrm{s}$,
forcing $\kN \to 0$ --- postictal silence~\cite{Perucca2014}.

\textbf{Phase 4 --- Hierarchical recovery.}
Microcircuits recover in seconds, areas in minutes,
whole-cortex function in tens of minutes, explaining the graded
recovery of cognitive functions after GTCS.

\subsection{Quantitative predictions}

\begin{enumerate}
\item \textbf{Tonic phase duration:} $10$--$15\,\mathrm{s}$,
  set by the time for cerebrovascular oxygen delivery to fail
  at $6.7\times$ demand. Literature:
  $10$--$20\,\mathrm{s}$~\cite{Perucca2014}.~\checkmark

\item \textbf{Tonic-clonic transition:}
  $\kN_{\mathrm{crit}} = 1 - \mathrm{CBF}_{\mathrm{rest}}/
  \mathrm{CBF}_{\max} \approx 0.71$, at which
  $\fUV \approx 5\,\mathrm{Hz}$; only delta oscillations
  remain.~\checkmark

\item \textbf{Ictal HFO signature:}
  Immediately before termination, the spectral exponent should
  steepen from $\alpha = -0.45$ (rest) to $\alpha \to -0.53$,
  consistent with the deep-layer (L5/L6) value measured in the
  companion paper~\cite{Santos2025a}, as projection neurons
  driving the seizure dominate the local field potential.
  \emph{Testable with Neuropixels during induced seizures.}

\item \textbf{Postictal scaling:}
  Duration $\propto$ seizure spatial extent $\times$ regional
  timescale ($\sim 800\,\mathrm{ms}$). Focal: minutes;
  generalized: tens of minutes.~\checkmark
\end{enumerate}

%% ════════════════════════════════════════════════════════════
%%  IV. SEMANTIC INERTIA
%% ════════════════════════════════════════════════════════════
\section{\label{sec:inertia}Semantic Inertia and the
         Thermodynamics of Belief Change}

\begin{table}[h]
\caption{\label{tab:h_scale}Semantic inertia $\hneu$ across
cortical scales.}
\begin{ruledtabular}
\begin{tabular}{lrrr}
Scale & $\taupop$ & $\hneu$ & Amplification \\
\hline
Microcircuit    &  19.3\,ms &
  $5.5\!\times\!10^{-14}$\,J$\cdot$s & $1\times$ \\
Cortical area   &  58.1\,ms &
  $6.8\!\times\!10^{-12}$\,J$\cdot$s & $124\times$ \\
Cortical region & 309.8\,ms &
  $7.2\!\times\!10^{-12}$\,J$\cdot$s & $132\times$ \\
Whole cortex    & 1936\,ms  &
  $1.5\!\times\!10^{-9}$\,J$\cdot$s  & $27{,}600\times$ \\
\end{tabular}
\end{ruledtabular}
\end{table}

The $27{,}600\times$ amplification (Fig.~\ref{fig:h_neural})
formalizes the thermodynamic resistance of distributed cortical
representations to perturbation. A whole-cortex representation
--- a deeply held belief, procedural memory, or perceptual
schema --- requires $27{,}600\times$ more energy to modify than
a local microcolumn pattern. This is not a metaphor but a direct
consequence of Eq.~(\ref{eq:h_asymp}) applied at the scale
$N_s = 10^5$.

This scale-dependent inertia provides a thermodynamic framework
for several clinical phenomena:

\textbf{Neuroplasticity windows:} Critical periods for sensory
plasticity~\cite{Hensch2005} correspond to developmental stages
when $\hneu$ at the area scale is transiently reduced, allowing
large-scale representations to be modified at lower energetic
cost.

\textbf{Ketamine vs.\ SSRIs:} Ketamine disrupts microcircuit
inertia ($\hneu^{(\mathrm{micro})}$) acutely, producing rapid
antidepressant effects; SSRIs gradually reduce area-scale inertia
through synaptic remodeling over weeks.

\textbf{Psychedelic-assisted therapy:} A transient global
reduction of $\hneu$ via 5-HT$_{2A}$ receptor agonism
temporarily lowers the thermodynamic barrier to updating
distributed representations, consistent with the broad-spectrum
effects observed in recent clinical trials.

\textbf{Post-ictal cognitive impairment:} Following GTCS,
whole-cortex $\hneu$ is depleted ($\kN \to 0$); hierarchical
recovery from local (seconds) to regional (hours) explains
the graded return of cognitive functions.

%% ════════════════════════════════════════════════════════════
%%  V. DISCUSSION
%% ════════════════════════════════════════════════════════════
\section{\label{sec:discussion}Discussion}

\subsection{Unification of neural timescales}

The single equation $f_n^{(s)} = 1/(n\tauISI^{(s)}\sqrt{\kN})$
predicts all major neural oscillation frequencies across four
orders of magnitude in timescale (milliseconds to seconds) and
five orders of magnitude in spatial scale (100 neurons to
$10^5$ cortical regions), using a single free parameter
($\kN = 0.15$) derived from metabolic first
principles~\cite{Attwell2001}. No fitting to oscillation data
was performed; the parameter is entirely determined by the
energy budget of cortical tissue.

\subsection{Epilepsy as a thermodynamic phase transition}

The cascade UV catastrophe model predicts a first-order-like
phase transition: normal activity ($\kN = 0.15$) $\to$ ictal
activity ($\kN \to 1$) $\to$ postictal silence ($\kN \to 0$),
with a well-defined critical point at $\kN_{\mathrm{crit}}
\approx 0.71$ (the tonic-clonic transition). This framework
suggests a mechanistic account of anti-epileptic drug action:
agents acting on sodium channels control $\tauabs$ and therefore
$\fUV^{\mathrm{neural}}$, raising the Planck cutoff and
preventing the cascade from initiating; agents increasing
inhibitory tone act by maintaining $\kN$ below the critical
threshold.

\subsection{Relation to cortical criticality}

The companion paper~\cite{Santos2025a} showed that MICrONS EM
data yields $\lambda_L = -0.082$ (near-critical). The epileptic
state ($\kN \to 1$, $\lambda_L > 0$) and the postictal state
($\kN \to 0$, $\lambda_L \ll -1$) represent the supercritical
and subcritical regimes flanking the normal operating point.
The tonic-clonic sequence is a thermodynamically driven
trajectory from near-critical through supercritical to
subcritical in the $(\kN, \lambda_L)$ phase space.

\subsection{Testable predictions and limitations}

The key empirical prediction --- $\alpha \to -0.53$ in ictal
LFP before termination --- is testable with Neuropixels during
pharmacologically induced seizures in rodents, or from existing
intracranial recordings in epilepsy monitoring units. The
postictal scaling prediction is testable from clinical EEG
databases with seizure annotations and neuropsychological
assessments.

The hierarchical $\tauISI^{(s)}$ values used ($150$, $800$,
$5{,}000\,\mathrm{ms}$) are estimates from the literature on
inter-areal communication delays and are not directly measured
within this framework. Future work should derive these from
electrocorticography during graded cortical stimulation across
anatomical distances. The fMRI prediction (modes
$f_n^{(\mathrm{ctx})}$ of the whole-cortex ensemble) requires
short-TR data ($\mathrm{TR} \leq 0.7\,\mathrm{s}$, band
$0.005$--$0.5\,\mathrm{Hz}$) for direct testing and is left as
a falsifiable prediction for future work with HCP 7T data.

%% ════════════════════════════════════════════════════════════
%%  VI. CONCLUSIONS
%% ════════════════════════════════════════════════════════════
\section{\label{sec:conclusions}Conclusions}

The metabolic constraint $\kN \approx 0.15$ is a scale-invariant
property of cortical tissue that quantizes neural oscillations
from gamma rhythms ($52\,\mathrm{Hz}$) at the microcircuit
scale to slow cortical rhythms ($0.5\,\mathrm{Hz}$) at the
regional scale, using a single equation with no free parameters
beyond the metabolic budget. Violation of this constraint
produces a metabolic UV catastrophe: the $6.7\times$ energy
amplification at full activation exceeds the cerebrovascular
reserve, generating the thermodynamic cascade of tonic activity,
clonic activity, and postictal silence that defines GTCS.
The semantic inertia amplifies $27{,}600\times$ from microcircuit
to whole cortex, providing a thermodynamic basis for the
resistance of distributed cortical representations to
perturbation, with direct implications for neuroplasticity,
psychopharmacology, and post-ictal recovery. Together with the
companion paper~\cite{Santos2025a}, these results constitute a
unified thermodynamic theory of cortical dynamics across all
scales.

%% ════════════════════════════════════════════════════════════
%%  ACKNOWLEDGMENTS
%% ════════════════════════════════════════════════════════════
\begin{acknowledgments}
The author thanks the Blue Brain Project and the Open Brain
Institute for publicly releasing the neocortical connectivity
data, and the MICrONS Consortium for the electron microscopy
reconstruction. The hierarchical scaling analysis was developed
as an extension of the spectral invariant results reported in
the companion paper~\cite{Santos2025a}.
\end{acknowledgments}

\section*{Data Availability}
Analysis code available at \texttt{github.com/colisor/TCSV}.
No new experimental data were collected for this paper.
All theoretical predictions are derived from the TCSV framework
and validated against published literature values.

%% ════════════════════════════════════════════════════════════
%%  FIGURES  (4 figures — fMRI section removed)
%% ════════════════════════════════════════════════════════════

\begin{figure}[t]
  \includegraphics[width=\columnwidth]{paper2_fig1_hierarquia}
  \caption{\label{fig:hierarchy}
    \textbf{Hierarchical scaling of cortical dynamics.}
    \textit{Left}: Frequency bands predicted by
    $f_n^{(s)} = 1/(n\tauISI^{(s)}\sqrt{\kN})$ at each
    cortical scale. Each bar spans from $f_6$ (slowest mode)
    to $f_1$ (fastest mode). The same equation with
    $\kN = 0.15$ invariant covers gamma ($\sim 52\,\mathrm{Hz}$)
    at the microcircuit scale to slow cortical oscillations
    ($\sim 0.5\,\mathrm{Hz}$) at the regional scale.
    \textit{Right}: Population integration time
    $\taupop = \tauISI^{(s)}\sqrt{\kN}$ by scale, spanning
    $19\,\mathrm{ms}$ (microcircuit) to
    $1{,}936\,\mathrm{ms}$ (whole cortex).}
\end{figure}

\begin{figure}[t]
  \includegraphics[width=\columnwidth]{paper2_fig2_epilepsia}
  \caption{\label{fig:epilepsy}
    \textbf{Epileptic seizure as metabolic UV catastrophe.}
    \textit{Upper left}: $\kN$ trajectory during a GTCS.
    Normal state ($\kN = 0.15$, dashed green); vascular
    threshold ($\kN = 0.71$, dotted orange). The four phases
    correspond to UV catastrophes at increasing spatial scales.
    \textit{Upper right}: Ensemble UV cutoff
    $\fUV^{\mathrm{ens}}$ vs.\ $\kN$; at $\kN = 0.9$ only
    delta activity ($\approx 2\,\mathrm{Hz}$) is sustainable.
    \textit{Lower left}: Cortical power $P(\kN)$ vs.\ $\kN$.
    Full activation ($80\,\mathrm{W}$) exceeds the
    cerebrovascular reserve ($42\,\mathrm{W}$); hypoxia
    develops within $10$--$15\,\mathrm{s}$.
    \textit{Lower right}: Quantitative predictions vs.\
    literature values (see Section~\ref{sec:epilepsy}).}
\end{figure}

\begin{figure}[t]
  \includegraphics[width=\columnwidth]{paper2_fig4_h_neural_escala}
  \caption{\label{fig:h_neural}
    \textbf{Semantic inertia $\hneu$ across cortical scales.}
    \textit{Left}: $\log_{10}(\hneu)$ by scale, from
    $5.5\times10^{-14}\,\mathrm{J\cdot s}$ (microcircuit,
    ``local belief'') to $1.5\times10^{-9}\,\mathrm{J\cdot s}$
    (whole cortex, ``cultural paradigm'') --- a $27{,}600\times$
    amplification reflecting the thermodynamic cost of modifying
    distributed vs.\ local representations.
    \textit{Right}: Clinical implications. The therapeutic
    timescale and mechanism of ketamine (acute, microcircuit),
    SSRIs (weeks, area scale), psychedelics (single session,
    whole cortex), and psychotherapy (months, regional scale)
    each correspond to the $\hneu$ target of that intervention.}
\end{figure}

\begin{figure}[t]
  \includegraphics[width=\columnwidth]{paper2_fig5_unificado}
  \caption{\label{fig:unified}
    \textbf{Unified TCSV framework across scales.}
    \textit{Center}: Four-level cortical hierarchy with
    $\kN = 0.15$ invariant (purple badge). Each level is
    characterized by its unit type, $\tauISI^{(s)}$, dominant
    frequency $f_1$, and semantic inertia $\hneu$.
    \textit{Top left}: $\kN$ trajectory during GTCS
    (summary of Fig.~\ref{fig:epilepsy}).
    \textit{Bottom left}: $\hneu$ by scale
    (summary of Fig.~\ref{fig:h_neural}).
    \textit{Bottom right}: Unified oscillation spectrum
    showing that a single equation governs cortical dynamics
    across five orders of magnitude in spatial scale.}
\end{figure}

%% ════════════════════════════════════════════════════════════
%%  BIBLIOGRAPHY
%% ════════════════════════════════════════════════════════════
\begin{thebibliography}{99}

\bibitem{Santos2025a}
G.~Santos,
\emph{Power-law spectral structure of synaptic connectivity
matrices is conserved across cortical areas, layers, and species},
Phys.\ Rev.\ E (submitted, 2025).

\bibitem{Attwell2001}
D.~Attwell and S.~B. Laughlin,
J.\ Cereb.\ Blood Flow Metab.\ \textbf{21}, 1133 (2001).

\bibitem{WHO2023}
World Health Organization,
\emph{Epilepsy Fact Sheet} (2023).
\url{https://www.who.int/news-room/fact-sheets/detail/epilepsy}

\bibitem{Stafstrom2010}
C.~E. Stafstrom and L.~Carmant,
Cold Spring Harb.\ Perspect.\ Med.\ \textbf{5}, a022533 (2015).

\bibitem{Buzsaki2007}
G.~Buzs\'aki and A.~Draguhn,
Science \textbf{304}, 1926 (2004).

\bibitem{Krishnan2011}
G.~P. Krishnan \emph{et al.},
J.\ Comput.\ Neurosci.\ \textbf{31}, 269 (2011).

\bibitem{Perucca2014}
P.~Perucca \emph{et al.},
Epilepsia \textbf{55}, 1065 (2014).

\bibitem{McClelland1995}
J.~L. McClelland \emph{et al.},
Psychol.\ Rev.\ \textbf{102}, 419 (1995).

\bibitem{Hensch2005}
T.~K. Hensch,
Nat.\ Rev.\ Neurosci.\ \textbf{6}, 877 (2005).

\bibitem{Beggs2003}
J.~M. Beggs and D.~Plenz,
J.\ Neurosci.\ \textbf{23}, 11167 (2003).

\bibitem{Mountcastle1997}
V.~B. Mountcastle,
Brain \textbf{120}, 701 (1997).

\end{thebibliography}

\end{document}